\chapter{绪论}
\label{chap:intro}

\section{研究背景与意义}
本文围绕课题“\ThesisTitle”展开。可信的实证结论要求每一个报告数值都能追溯到实际执行的实验，
每一条引用都能对应到可核验的出版物。因此，本研究以一条自动化、可审计的科研流程为核心组织全部工作。

\section{研究目标}
\begin{itemize}
  \item 建立完全由声明式课题配置驱动、可复现的实验方案；
  \item 在\MetricPrimaryName{}及各次要指标上，以满足前提假设的统计检验量化本文配置相对其他方案的表现；
  \item 正文只报告实验生成的数值，并以溯源哈希将论文与结果包相互绑定。
\end{itemize}

\section{论文结构}
第~\ref{chap:related}~章综述经 DOI 核验的相关文献；第~\ref{chap:method}~章介绍研究方法与统计方案；
第~\ref{chap:experiments}~章报告实验结果；第~\ref{chap:conclusion}~章总结全文。
